\scanpage{8}
% 1–4
\deleted{\uline{Cor} $A_1$ est stable par primitives canoniques (i.d. prises d'augmentation nulle) contient $\log(1+t_1)$\unsure{,|;} (primitive de $\frac{1}{1+t_1}$) donc $\exp.t_1$ (série inverse), les fonctions $\cos t_1$, $\sin t_1$, $\mathrm{tg}\,t_1$ sont $\in A_1$ .....}

% 5–6
\deleted{\uline{Lemme} \deleted{\unsure{h)|g)|1)} Soit} Soit $X$ un \inserted{$A$-}schéma formel (lisse) \inserted{de dim $n=p+q$}, $D$ un ``système de Pfaff de degré $p$ complètement intégrable sur $X$'', \deleted{\unsure{D}}}

% 7–10
\deleted{\deleted{deux \unsure{soient|lisses} Soient $Y,Z$ des \inserted{$A$-}schémas formels \inserted{(lisses)} de dims $p$ resp $q$,} $p:X\to Y$ et $i:Z\to X$ \inserted{des $A$-morphismes}, tels que $T(p)(0)$ soit un \deleted{\unsure{appl.}}\inserted{isom} $D(0)-T(Y)_0$ et que \deleted{$T(i)(0)$} le composé \deleted{$T(i)(0)$}}
\[
 \deleted{T(Z)_{0}\deleted{(0)}\xrightarrow{T(i)_0}T(X)_{0}\deleted{(0)}\xrightarrow{\mathrm{can}}T(X)_0/D(0)}
\]

% 11–14
\deleted{soit un isom. On sait qu'il existe alors un isomorphisme \inserted{formel (pas néc. $A$-isom)} unique}
\[
 \deleted{X\overset{\varphi}{\simeq}Y\times Z}
\]
\deleted{\deleted{\unsure{transformant|transportant} et \unsure{qui|puis} tel que $\mathrm{pr}_1\circ\varphi=p$, $\varphi\circ i=(z\mapsto(0,z))$} et}

% 15–19
\deleted{Soit $Z$ \inserted{$A$-}schéma formel \inserted{(lisse)} de dim $q$, $i:Z-X$ une $A$-\unsure{immersion|réimmersion} telle que le composé}
\[
 \deleted{T(Z)_0\xrightarrow{T(i)_0}T(X)_0\longrightarrow T(X)_0/D(0)}
\]
\deleted{soit un isom. On sait qu'il existe alors un unique morphisme formel}
\[
 \deleted{X\xrightarrow{\pi}Z}
\]

% 20–25
\deleted{tel que $\pi\circ i=\mathrm{id}$, et que ``$\pi$ soit \deleted{\unsure{constante sur|constant sur|constant sur les}}\inserted{compatible avec} le feuilletage défini par $D$''. Conditions équivalentes}

\deleted{(i) $\pi$ est un $A$-morphisme}

\deleted{(ii) $\exists$ $A$-isom $X\simeq B\times Y$, tel que le feuilletage formel de $X$ défini par $D$ soit le feuilletage défini par la projection $\mathrm{pr}_1:X\to B$. On dit}
% archive sheet number: 7
